\documentclass[../../main.tex]{subfiles}

\begin{document}

\section{导数的积分: 微分不定积分}\label{section:ch06-导数积分}

令 $f$ 为闭有界区间 $[a,b]$ 上的连续函数. 在等式\eqref{align:ch06-换元}中, 取 $a=u$, $b=v$ 得到以下离散形式的计算积分的基本定理:
\begin{align*}
\int_{a}^{b}\operatorname{Diff}_{h}f=\operatorname{Avg}_{h}f(b)-\operatorname{Avg}_{h}f(a).
\end{align*}
由于 $f$ 是连续的, 当 $h\to 0^{+}$ 时右边的极限等于 $f(b)-f(a)$. 我们现在证明: 若 $f$ 是绝对连续的, 则当 $h\to 0^{+}$ 时左边的极限等于 $\int_{a}^{b}f'$, 从而对 Lebesgue 积分证明了计算积分的基本定理.

\begin{theorem}\label{theorem:ch06-AC导数积分}
令函数 $f$ 在闭有界区间 $[a,b]$ 上是绝对连续的. 则 $f$ 在 $(a,b)$ 上是几乎处处可微的, 它的导数 $f'$ 在 $[a,b]$ 上是可积的, 且
\begin{align}
\int_{a}^{b}f'=f(b)-f(a).\label{align:ch06-积分基本等式}
\end{align}
\end{theorem}

\begin{proof}
我们从离散形式的积分计算的基本定理推出
\begin{align}
\lim_{n\to\infty}\left[\int_{a}^{b}\operatorname{Diff}_{1/n}f\right]=f(b)-f(a).\label{align:ch06-离散极限}
\end{align}
定理\cref{theorem:ch06-AC是递增差}告诉我们 $f$ 是 $[a,b]$ 上的递增函数的差, 因此, 根据 Lebesgue 定理, 在 $(a,b)$ 上是几乎处处可微的. 因此在 $(a,b)$ 上 $\{\operatorname{Diff}_{1/n}f\}$ 几乎处处逐点收敛于 $f'$. 另一方面, 根据定理\cref{theorem:ch06-AC与一致可积}, $\{\operatorname{Diff}_{1/n}f\}$ 在 $[a,b]$ 上是一致可积的. Vitali 收敛定理 (见 4.6 节) 允许在积分号下取极限, 得出
\begin{align}
\lim_{n\to\infty}\left[\int_{a}^{b}\operatorname{Diff}_{1/n}f\right]=\int_{a}^{b}\lim_{n\to\infty}\operatorname{Diff}_{1/n}f=\int_{a}^{b}f'.\label{align:ch06-Vitali极限}
\end{align}
公式\eqref{align:ch06-积分基本等式}从\eqref{align:ch06-离散极限}和\eqref{align:ch06-Vitali极限}得出.
\end{proof}

在微积分的研究中, 不定积分是对 Riemann 积分定义的. 我们这里称闭有界区间 $[a,b]$ 上的函数 $f$ 为 $g$ 在 $[a,b]$ 上的不定积分, 若 $g$ 在 $[a,b]$ 上是 Lebesgue 可积的且
\begin{align}
\text{对所有 }x\in[a,b],\quad f(x)=f(a)+\int_{a}^{x}g.\label{align:ch06-不定积分}
\end{align}

\begin{theorem}\label{theorem:ch06-AC与不定积分}
闭有界区间 $[a,b]$ 上的函数 $f$ 在 $[a,b]$ 上是绝对连续的当且仅当它是 $[a,b]$ 上的不定积分.
\end{theorem}

\begin{proof}
首先假定 $f$ 在 $[a,b]$ 上是绝对连续的. 对每个 $x\in(a,b)$, $f$ 在 $[a,x]$ 上是绝对连续的, 因此, 根据前一个定理, 在 $[a,b]$ 换为 $[a,x]$ 的情形,
\begin{align*}
f(x)=f(a)+\int_{a}^{x}f'.
\end{align*}
因此 $f$ 是 $f'$ 在 $[a,b]$ 上的不定积分. 反之, 假定 $f$ 是 $g$ 在 $[a,b]$ 上的不定积分. 对于 $(a,b)$ 中的开区间的不交族 $\{(a_{k},b_{k})\}_{k=1}^{n}$, 若我们定义 $E=\bigcup_{k=1}^{n}(a_{k},b_{k})$, 则根据积分的单调性和在区域上的可加性,
\begin{align}
\sum_{k=1}^{n}|f(b_{k})-f(a_{k})|=\sum_{k=1}^{n}\left|\int_{b_{k}}^{a_{k}}g\right|\leqslant\sum_{k=1}^{n}\int_{b_{k}}^{a_{k}}|g|=\int_{E}|g|.\label{align:ch06-绝对连续估计}
\end{align}
令 $\varepsilon>0$. 由于 $|g|$ 在 $[a,b]$ 上是可积的, 根据第 4 章的命题 23, 存在 $\delta>0$, 使得若 $E\subseteq[a,b]$ 是可测的且 $m(E)<\delta$, 则 $\int_{E}|g|<\varepsilon$. 从\eqref{align:ch06-绝对连续估计}得出同样的 $\delta$ 对应于 $f$ 在 $[a,b]$ 上绝对连续准则的 $\varepsilon$ 挑战.
\end{proof}

\begin{corollary}\label{corollary:ch06-单调绝对连续}
令函数 $f$ 在闭有界区间 $[a,b]$ 上单调. 则 $f$ 在 $[a,b]$ 上是绝对连续的当且仅当
\begin{align}
\int_{a}^{b}f'=f(b)-f(a).\label{align:ch06-单调等式}
\end{align}
\end{corollary}

\begin{proof}
定理\cref{theorem:ch06-AC导数积分}断言若 $f$ 是绝对连续的则\eqref{align:ch06-单调等式}成立, 不用任何单调性假设. 反之, 假设 $f$ 是递增的且\eqref{align:ch06-单调等式}成立. 令 $x$ 属于 $[a,b]$. 根据积分在区域上的可加性,
\begin{align*}
0=\int_{a}^{b}f'-[f(b)-f(a)]=\left\{\int_{a}^{x}f'-[f(x)-f(a)]\right\}+\left\{\int_{x}^{b}f'-[f(b)-f(x)]\right\}.
\end{align*}
根据推论\cref{corollary:ch06-单调导数可积},
\begin{align*}
\int_{a}^{x}f'-[f(x)-f(a)]\leqslant 0\quad\text{且}\quad\int_{x}^{b}f'-[f(b)-f(x)]\leqslant 0.
\end{align*}
若两个非负数的和为零, 则它们都为零. 因此,
\begin{align*}
f(x)=f(a)+\int_{a}^{x}f'.
\end{align*}
因此 $f$ 是 $f'$ 的不定积分. 前一个定理告诉我们 $f$ 是绝对连续的.
\end{proof}

\begin{lemma}\label{lemma:ch06-积分为零}
令 $f$ 在闭有界区间 $[a,b]$ 上是可积的. 则
\begin{align}
\text{对几乎所有 }x\in[a,b],\quad f(x)=0\label{align:ch06-几乎处处零}
\end{align}
当且仅当对所有 $(x_{1},x_{2})\subseteq[a,b]$,
\begin{align}
\int_{x_{1}}^{x_{2}}f=0.\label{align:ch06-区间积分为零}
\end{align}
\end{lemma}

\begin{proof}
显然\eqref{align:ch06-几乎处处零}蕴涵\eqref{align:ch06-区间积分为零}. 反之, 假定\eqref{align:ch06-区间积分为零}成立. 我们宣称
\begin{align}
\text{对所有可测子集 }E\subseteq[a,b],\quad\int_{E}f=0.\label{align:ch06-任意集积分为零}
\end{align}
事实上, \eqref{align:ch06-任意集积分为零}对所有包含于 $(a,b)$ 的开集成立, 这是由于积分是可数可加的且每个开集是开区间的可数不交族的并. 积分的连续性告诉我们\eqref{align:ch06-任意集积分为零}对所有包含于 $(a,b)$ 的 $G_{\delta}$ 集也成立, 这是由于每个这样的集是开集的可数下降族的交. 但 $[a,b]$ 的每个可测子集形如 $G\sim E_{0}$, 其中 $G$ 是 $(a,b)$ 的 $G_{\delta}$ 子集而 $m(E_{0})=0$ (见 2.4 节定理 11). 我们从积分在区域上的可加性得出\eqref{align:ch06-任意集积分为零}. 定义
\begin{align*}
E^{+}=\{x\in[a,b]\mid f(x)\geqslant 0\}\quad\text{与}\quad E^{-}=\{x\in[a,b]\mid f(x)\leqslant 0\}.
\end{align*}
它们是 $[a,b]$ 的两个可测子集, 因此, 根据\eqref{align:ch06-任意集积分为零},
\begin{align*}
\int_{a}^{b}f^{+}=\int_{E^{+}}f=0\quad\text{且}\quad\int_{a}^{b}(-f^{-})=-\int_{E^{-}}f=0.
\end{align*}
根据第 4 章的命题 9, 积分为零的非负可积函数必须在它的定义域上几乎处处消失. 于是 $f^{+}$ 和 $f^{-}$ 在 $[a,b]$ 上几乎处处消失, 因此 $f$ 也如此.
\end{proof}

\begin{theorem}\label{theorem:ch06-微分基本定理}
令 $f$ 为闭有界区间 $[a,b]$ 上的可积函数. 则对几乎所有 $x\in(a,b)$,
\begin{align}
\frac{d}{dx}\left[\int_{a}^{x}f\right]=f(x).\label{align:ch06-微分等式}
\end{align}
\end{theorem}

\begin{proof}
对所有 $x\in[a,b]$, 定义 $[a,b]$ 上的函数 $F$ 为 $F(x)=\int_{a}^{x}f$. 定理\cref{theorem:ch06-AC与不定积分}告诉我们由于 $F$ 是不定积分, 它是绝对连续的. 因此, 根据定理\cref{theorem:ch06-AC导数积分}, $F$ 在 $(a,b)$ 上是几乎处处可微的且它的导数 $F'$ 是可积的. 根据前一个引理, 为证明可积函数 $F'-f$ 在 $[a,b]$ 上几乎处处消失, 仅须证明它在 $[a,b]$ 的每个闭子区间上的积分是零. 令 $[x_{1},x_{2}]$ 包含于 $[a,b]$. 根据定理\cref{theorem:ch06-AC导数积分} (在 $[a,b]$ 被 $[x_{1},x_{2}]$ 替代的情形), 以及积分的线性与在区域上的可加性,
\begin{align*}
\int_{x_{1}}^{x_{2}}[F'-f]=\int_{x_{1}}^{x_{2}}F'-\int_{x_{1}}^{x_{2}}f=F(x_{2})-F(x_{1})-\int_{x_{1}}^{x_{2}}f=\int_{a}^{x_{2}}f-\int_{a}^{x_{1}}f-\int_{x_{1}}^{x_{2}}f=0.
\end{align*}
\end{proof}

一个有界变差函数称为是\textbf{奇异的}, 若它的导数几乎处处消失. Cantor-Lebesgue 函数是一个非常数的奇异函数. 我们从定理\cref{theorem:ch06-AC导数积分}推出一个绝对连续的函数是奇异的当且仅当它是常数. 令 $f$ 为 $[a,b]$ 上的有界变差函数. 根据推论\cref{corollary:ch06-BV可微}, $f'$ 在 $[a,b]$ 上是可积的. 对所有 $x\in[a,b]$, 定义
\begin{align*}
g(x)=\int_{a}^{x}f'\quad\text{与}\quad h(x)=f(x)-\int_{a}^{x}f',
\end{align*}
使得在 $[a,b]$ 上 $f=g+h$.

根据定理\cref{theorem:ch06-AC与不定积分}, 函数 $g$ 是绝对连续的. 我们从定理\cref{theorem:ch06-微分基本定理}推出函数 $h$ 是奇异的. 以上将有界变差函数 $f$ 表示为有界变差函数 $g$ 和 $h$ 的和 $g+h$ 的分解, 其中 $g$ 是绝对连续的而 $h$ 是奇异的, 称为 $f$ 的 Lebesgue 分解.

\begin{exercise}\label{exercise:ch06-48}
Cantor-Lebesgue 函数 $\varphi$ 在 $[0,1]$ 上是连续与递增的. 从定理\cref{theorem:ch06-AC导数积分}得出 $\varphi$ 在 $[0,1]$ 上不是绝对连续的. 比较这个推理与习题 40 提出的推理.
\end{exercise}

\begin{exercise}\label{exercise:ch06-49}
令 $f$ 为 $[a,b]$ 上连续且在 $(a,b)$ 上几乎处处可微的函数. 证明
\begin{align*}
\int_{a}^{b}f'=f(b)-f(a)
\end{align*}
当且仅当
\begin{align*}
\int_{a}^{b}\left[\lim_{n\to\infty}\operatorname{Diff}_{1/n}f\right]=\lim_{n\to\infty}\left[\int_{a}^{b}\operatorname{Diff}_{1/n}f\right].
\end{align*}
\end{exercise}

\begin{exercise}\label{exercise:ch06-50}
令 $f$ 为 $[a,b]$ 上连续且在 $(a,b)$ 上几乎处处可微的函数. 证明: 若 $\{\operatorname{Diff}_{1/n}f\}$ 在 $[a,b]$ 上一致可积, 则
\begin{align*}
\int_{a}^{b}f'=f(b)-f(a).
\end{align*}
\end{exercise}

\begin{exercise}\label{exercise:ch06-51}
令 $f$ 为 $[a,b]$ 上连续且在 $(a,b)$ 上几乎处处可微的函数. 假定存在非负函数 $g$, 该函数在 $[a,b]$ 上可积且对所有 $n$, 在 $[a,b]$ 上 a.e. $|\operatorname{Diff}_{1/n}f|\leqslant g$. 证明
\begin{align*}
\int_{a}^{b}f'=f(b)-f(a).
\end{align*}
\end{exercise}

\begin{exercise}\label{exercise:ch06-52}
令 $f$ 和 $g$ 在 $[a,b]$ 上绝对连续. 证明
\begin{align*}
\int_{a}^{b}f\cdot g'=f(b)g(b)-f(a)g(a)-\int_{a}^{b}f'\cdot g.
\end{align*}
\end{exercise}

\begin{exercise}\label{exercise:ch06-53}
令函数 $f$ 在 $[a,b]$ 上绝对连续. 证明 $f$ 在 $[a,b]$ 上是 Lipschitz 的当且仅当存在 $c>0$, 使得在 $[a,b]$ 上 a.e. $|f'|\leqslant c$.
\end{exercise}

\begin{exercise}\label{exercise:ch06-54}
(i) 令 $f$ 为 $[a,b]$ 上的奇异递增函数. 用 Vitali 覆盖引理证明 $f$ 具有以下性质: 给定 $\varepsilon>0$, $\delta>0$, 存在 $(a,b)$ 中的开区间的有限不交族 $\{(a_{k},b_{k})\}_{k=1}^{n}$ 使得
\begin{align*}
\sum_{k=1}^{n}[b_{k}-a_{k}]<\delta\quad\text{且}\quad\sum_{k=1}^{n}[f(b_{k})-f(a_{k})]>f(b)-f(a)-\varepsilon.
\end{align*}

(ii) 令 $f$ 为具有 (i) 部分描述的性质的 $[a,b]$ 上的递增函数. 证明 $f$ 是奇异的.

(iii) 令 $\{f_{n}\}$ 为 $[a,b]$ 上的奇异递增函数序列, 使得对每个 $x\in[a,b]$, 级数 $\sum_{n=1}^{\infty}f_{n}(x)$ 收敛到有限值. 对 $x\in[a,b]$, 定义
\begin{align*}
f(x)=\sum_{n=1}^{\infty}f_{n}(x).
\end{align*}
证明 $f$ 也是奇异的.
\end{exercise}

\begin{exercise}\label{exercise:ch06-55}
令 $f$ 为 $[a,b]$ 上的有界变差函数, 对所有 $x\in[a,b]$, 定义 $v(x)=TV(f_{[a,x]})$.
\begin{enumerate}[(i)]
\item 证明在 $[a,b]$ 上 a.e. $|f'|\leqslant v'$, 且由此推出
\begin{align*}
\int_{a}^{b}|f'|\leqslant TV(f).
\end{align*}
\item 证明以上不等式成为等式当且仅当 $f$ 在 $[a,b]$ 上是绝对连续的.
\item 将 (i) 部分和 (ii) 部分分别与推论\cref{corollary:ch06-单调导数可积}和推论\cref{corollary:ch06-单调绝对连续}比较.
\end{enumerate}
\end{exercise}

\begin{exercise}\label{exercise:ch06-56}
令 $g$ 在 $[a,b]$ 上严格递增且绝对连续.
\begin{enumerate}[(i)]
\item 证明对于 $(a,b)$ 的任何开子集 $\mathcal{O}$,
\begin{align*}
m(g(\mathcal{O}))=\int_{\mathcal{O}}g'(x)\,dx.
\end{align*}
\item 证明对于 $(a,b)$ 的任何 $G_{\delta}$ 子集 $E$,
\begin{align*}
m(g(E))=\int_{E}g'(x)\,dx.
\end{align*}
\item 证明对于 $[a,b]$ 的任何测度为零的子集 $E$, 它的象 $g(E)$ 也有零测度, 因此
\begin{align*}
m(g(E))=0=\int_{E}g'(x)\,dx.
\end{align*}
\item 证明对于 $[a,b]$ 的任何可测子集 $A$,
\begin{align*}
m(g(A))=\int_{A}g'(x)\,dx.
\end{align*}
\item 令 $c=g(a)$ 而 $d=g(b)$. 证明对于 $[c,d]$ 上的任何简单函数 $\varphi$,
\begin{align*}
\int_{c}^{d}\varphi(y)\,dy=\int_{a}^{b}\varphi(g(x))g'(x)\,dx.
\end{align*}
\item 证明对于任何 $[c,d]$ 上的非负可积函数 $f$,
\begin{align*}
\int_{c}^{d}f(y)\,dy=\int_{a}^{b}f(g(x))g'(x)\,dx.
\end{align*}
\item 证明在 $f$ 是 $g(\mathcal{O})$ 的特征函数与复合有定义的情形下, (i) 部分可从 (vi) 部分得到.
\end{enumerate}
\end{exercise}

\begin{exercise}\label{exercise:ch06-57}
在前一个习题的 (vi) 部分, 若我们仅假设 $g$ 是递增的, 不必是严格的, 换元公式是否成立?
\end{exercise}

\begin{exercise}\label{exercise:ch06-58}
构造一个 $[0,1]$ 上绝对连续的严格递增函数 $f$, 使得在一个正测度集上 $f'=0$. (提示: 令 $E$ 为具有正测度的广义 Cantor 集在 $[0,1]$ 的相对补集而 $f$ 是 $\chi_{E}$ 的不定积分. 这样的 Cantor 集的构造见第 2 章的习题 39.)
\end{exercise}

\begin{exercise}\label{exercise:ch06-59}
对于 $[c,d]$ 上的非负可积函数 $f$ 和 $[a,b]$ 上使得 $g([a,b])\subseteq[c,d]$ 的严格递增的绝对连续函数 $g$, 是否可能通过证明对几乎所有 $x\in(a,b)$,
\begin{align*}
\frac{d}{dx}\left[\int_{g(a)}^{g(x)}f(s)\,ds-\int_{a}^{x}f(g(t))g'(t)\,dt\right]=0
\end{align*}
来证明换元公式
\begin{align*}
\int_{g(a)}^{g(b)}f(y)\,dy=\int_{a}^{b}f(g(x))g'(x)\,dx
\end{align*}
的合理性?
\end{exercise}

\begin{exercise}\label{exercise:ch06-60}
令 $f$ 在 $[a,b]$ 上绝对连续与奇异. 证明 $f$ 是常数. 也证明一个有界变差函数的 Lebesgue 分解是唯一的, 若奇异函数要求在 $x=a$ 消失.
\end{exercise}

\end{document}